\section{Good-orbit estimates for the actual biased selection}
\label{sec:biased-orbits}

\status{This proves the nonconstant Bernoulli estimate needed by the
arithmetic top-set construction. It allows an arbitrary inclusion profile
between $3\delta/4$ and $\delta$, keeps the actual parity and forced-point
conditions, and supplies the weighted expectation used in the subsequent
finite selection. It does not construct an arithmetic host.}

Fix $\epsilon>0$ and real constants $C_o,L_b$. Put
\[
\begin{gathered}
 \theta=\epsilon/8,\quad h=2+8\epsilon,\quad
 K=\lceil(h+2)/\theta\rceil_{\mathbb N},\quad r=K+1,
 \quad\delta=x^{-\theta},\\
 u=\lceil\delta^{-1}\rceil_{\mathbb N},\quad
 \ell=(6r+1)u,\quad\pi_0=10^{-(6r+1)},\\
 M=\lceil(h+3)\log x/\pi_0\rceil_{\mathbb N},\quad
 G=(M+5)\ell,\\
 N=\left\lceil\max\{6eG\delta,(h+3)\log_2x\}\right\rceil_{\mathbb N}+1.
\end{gathered}                                                   \tag{BO.1}
\]
These are the literal \texttt{KE}, \texttt{rE}, \texttt{goodU},
\texttt{orbitSmallMultiplicity}, \texttt{goodG'} and
\texttt{orbitWindowCount}. Choose each coordinate independently, with
probability $q_i\in[3\delta/4,\delta]$.

\paragraph{BO1: the finite product law and count tails.}
On a finite set $S$, a subset $U$ has weight
$\prod_{i\in U}q_i\prod_{i\in S\setminus U}(1-q_i)$.
Summing over any unused coordinates multiplies by
$\prod_i(q_i+1-q_i)=1$. Consequently conditioning on a selected
coordinate $i_0$ leaves the other coordinates independent with their
original probabilities. Events on disjoint sets still factor.

For count tails, compare each coordinate with a Bernoulli($\delta$)
coordinate using joint states $(0,0),(0,1),(1,1)$ of probabilities
$1-\delta,\delta-q_i,q_i$. Padding to $g$ coordinates gives, for
$|S|\le g$,
\[
 \Pr(|U|\ge b)\le {g\choose b}\delta^b,
 \qquad
 \Pr(|U|\ge b)\le
 B(g,\delta,b):=\sum_{j=b}^{g}{g\choose j}\delta^j(1-\delta)^{g-j}.
                                                               \tag{BO.2}
\]
The first bound also follows by counting selected $b$-subsets. This
comparison is used only for increasing count events, not for marked tiles.

\paragraph{BO2: the biased marked-tile probability.}
A complete tile has $6r+1$ consecutive cells of length $u$. Its
odd-numbered cells must be empty; every even-numbered cell must contain
exactly one selected point, whose \emph{coordinate} is even.
Any such cell contains at least $\lfloor u/2\rfloor$ even coordinates.
The corresponding event probabilities are therefore at least
\[
 A_0=(1-\delta)^u,\qquad
 A_1=\lfloor u/2\rfloor\frac{3\delta}{4}(1-\delta)^{u-1}.
\]
The singleton events are disjoint, and all cells are disjoint, so the
actual marked probability is at least the source's
\[
 \pi=A_0^{3r+1}A_1^{3r}=\texttt{markedLow}(\delta,r,u).
                                                               \tag{BO.3}
\]
If $0<\delta\le1/100$, then $1\le u\delta\le1+\delta$ and
$\log(1-\delta)\ge-\delta/(1-\delta)$ give
$A_0\ge e^{-103/100}>3/10$. Also
$\lfloor u/2\rfloor\delta\ge99/200$, whence
$A_1\ge891/8000>1/10$. Thus
\[
 0\le\pi_0\le\pi\le1.                                        \tag{BO.4}
\]
This is the additional calculation needed for the biased law.

\paragraph{BO3: all three failures, with one point forced.}
Use the source's cyclic windows
$W(i,L)=\{(i+j)\bmod n:0\le j<L\}$. A set is good when all
$L_2$-windows contain at most $K$ points, every $M+1$ successive
complete tiles contains a marked tile, and all $G$-windows contain at
most $N$ points. Complete tile indices are read modulo
$t=\lfloor n/\ell\rfloor$.

Assume $(M+1)\ell\le n$ and force $i_0<n$ to be selected. A failed
short window still contains at least $K$ unforced selected points.
For each starting tile, its $M+1$ tiles are distinct; remove the at most
one containing $i_0$. At least $M$ independent, unaffected tiles remain.
If $i_0$ is in the incomplete remainder, none is removed. A failed
long window implies at least $N-1$ unforced selected points. Union bounds
over the $n$ starts, using $t\le n$, prove the exact source bound
\[
 \Pr(\mathrm{bad}\mid i_0\text{ selected})\le
 n{L_2\choose K}\delta^K+n(1-\pi)^M+nB(G,\delta,N-1).
                                                               \tag{BO.5}
\]
The long-window threshold is deliberately weakened to the source's
$N-1$. No cyclic rotation of the parity condition is used.

\paragraph{BO4: one explicit uniform threshold.}
Set
\[
\begin{gathered}
 C_1=(h+3)/\pi_0+2,\quad C_*=2(6r+1)C_1,\\
 x_0=\exp\max\{1,\,4C_*/\theta^2,\,\log(100)/\theta\},\qquad
 C_b=(|C_o|+1)\bigl((\max\{L_b,1\})^K+2\bigr).
\end{gathered}                                                   \tag{BO.6}
\]
For $x\ge x_0$, writing $y=\log x$ gives
$C_*(y+1)\le2C_*y\le(\theta y)^2/2\le e^{\theta y}$.
The exact ceiling bounds now yield
\[
 (M+1)\ell\le C_*(\log x+1)x^\theta
             \le x^{2\theta}.                                  \tag{BO.7}
\]
Hence the tile condition holds whenever
$x^{\epsilon/4}\le n\le C_ox^h$. If $L_2\le L_b$, the three terms
in (BO.5) are respectively at most
\[
 (|C_o|+1)(\max\{L_b,1\})^Kx^{-2},\qquad
 (|C_o|+1)x^{-3},\qquad (|C_o|+1)x^{-3}.                         \tag{BO.8}
\]
Indeed $\theta K\ge h+2$ and $\pi_0M\ge(h+3)\log x$.
For the last term put $b=N-1$. The factorial estimate
$b!\ge(b/e)^b$ and (BO.2) give
$B(G,\delta,b)\le6^{-b}\le e^{-b}$, since $b\ge6eG\delta$;
also $b\ge(h+3)\log_2x\ge(h+3)\log x$. The empty tail is harmless.
Consequently the conditional probability is at most $C_bx^{-2}$,
uniformly over every allowed length, profile, base and forced point.

\paragraph{BO5: the literal ambient-orbit and weighted statement.}
Let $A$ be a finite abelian group and choose its points independently
with probabilities $q(a)\in[3\delta/4,\delta]$. Fix a step $v$ of
order $n$, a point $a$, and any base $b$ with $a=b+i_0v$.
The map $i\mapsto b+iv$, $0\le i<n$, is injective. Marginalizing
outside its image and then factoring out $q(a)$ proves exactly
\[
 \Pr(a\in T\text{ and its based orbit is bad})
 =q(a)\Pr(\mathrm{bad}\mid i_0\text{ selected})
 \le\delta C_bx^{-2}.                                          \tag{BO.9}
\]
The base may depend on both the step and $a$, as the source's
$\mathrm{bs}(v,a)$ does. For a finite set of steps $v_p$ satisfying
the length bounds and arbitrary nonnegative weights $\omega_p$,
linearity gives
\[
 \mathbb E\sum_p\omega_p
   \#\{a\in T:\text{the prescribed based orbit is bad}\}
 \le \delta |A|C_bx^{-2}\sum_p\omega_p.                         \tag{BO.10}
\]
No independence between the different bad-orbit events is required.
This supplies the probability input to the simultaneous coset selection;
that selection and arithmetic realization are separate steps.
